\begin{theorem}[Proved: ordinary ultraproducts lose modular exactness]
\label{C4:ultraproduct-exclusion}
Let $\mathcal U$ be a nonprincipal ultrafilter on $\mathbb N$.
For arbitrary integers $r_i\geq2$, the ordinary category
ultraproduct
\[
 \mathcal C=\prod_{i\to\mathcal U}\mathcal A_{r_i}(\mathbb Q)
\]
is not exact over any nonzero unital ring and is not total. The same is
true of the ultrapower $\prod_{i\to\mathcal U}\Delta_{\rm inj}$,
although $\Delta_{\rm inj}$ itself is integrally exact and total.
Thus there is no bound on the support size of affine
equalizers in terms of the number of parallel arrows alone, even
for two arrows in $\Delta_{\rm inj}$.
\end{theorem}
\begin{proof}
We first give a uniform lower bound for one finite ordered-set
family. Let $P=[n]$, $Q=[n+1]$, and let
\[
 d_0(j)=j,\qquad d_1(j)=j+1 \quad (1\leq j\leq n).
 \tag{UP.1}
\]
For a finite ordered target $D$, identify an embedding
$P\to D$ with its increasing $n$-tuple of images. An embedding
$g:Q\to D$ gives an edge between its first and last $n$-tuples.
These edges form a simple undirected graph: an edge determines
the increasing $(n+1)$-tuple uniquely. Orient every edge from
the first window to the last window.

This graph has no odd closed walk of length $L<n$. Indeed,
write its successive $n$-tuples as $x^0,\ldots,x^L=x^0$.
For each forward step put $\epsilon_j=1$ and for each backward
step put $\epsilon_j=-1$, and set
$S_0=0$, $S_j=\epsilon_1+\cdots+\epsilon_j$. Since
$\max_j S_j-\min_j S_j\leq L<n$, there is an integer $h$ with
\[
 1+\max_j S_j\leq h\leq n+\min_j S_j.
\]
The coordinate at position $h-S_j$ survives every window shift;
thus $x^j_{h-S_j}=x^0_h$ for every $j$. At the end the distinct
entries of $x^L=x^0$ force $h-S_L=h$, hence $S_L=0$.
An odd number of signs $\pm1$ cannot have sum zero, proving the
claim. Only this unbounded lower bound is needed here.

In fact every closed walk of length $L<n$ has signed forward
minus backward edge count zero, by the same coordinate argument.
Let $R$ be any nonzero unital ring and let
$\mu\in R[\operatorname{Emb}(Q,D)]$ be an affine equalizer.
Regard its support as a finite oriented graph with edge coefficients
in $R$. The equality $\mu d_0=\mu d_1$ says that this chain is a
circulation. Choose a spanning forest of its support. Subtracting
multiples of the fundamental cycles removes each edge outside the
forest, and a circulation supported on a finite forest is zero,
by successively removing its leaves. If the support had fewer
than $n$ edges, each fundamental cycle would have length less
than $n$ and signed augmentation zero. Every circulation on that
support would then have augmentation zero, contrary to
$\operatorname{aug}\mu=1$. Hence
\[
 |\operatorname{supp}\mu|\geq n.
 \tag{UP.2}\label{C4:window-support}
\]
This argument uses only integer signs, so also works for
noncommutative unital rings. For $R=\mathbb F_2$ one can
alternatively note that an odd Eulerian edge set contains an odd
cycle. The same lower bound applies when $P,Q$ carry zero
cocycles in $\mathcal A_r(\mathbb Q)$ and $D$ carries any
cocycle: forgetting labels sends an affine equalizer to one for
(UP.1), preserves augmentation, and can only decrease support.

Take any sequence $n_i\to\infty$ and form the two arrows in
$\mathcal C$ represented by (UP.1) between the zero-cocycle
objects $[n_i]$ and $[n_i+1]$. Suppose they had an affine
equalizer with a finite number $m$ of nonzero terms. Represent
its target and its $m$ arrows by sequences in the factors, keeping their coefficients in $R$ fixed. The
formal equality of the two finite sums of composite arrows is
a finite collection of equalities, after grouping equal terms.
Each of these equalities holds on a member of $\mathcal U$;
their finite intersection is still a member of $\mathcal U$.
On that set of indices the representatives consequently give
an $R$-affine equalizer with at most $m$ terms and
augmentation one. But $n_i>m$ also holds on a member of
$\mathcal U$, contradicting (UP.2).

Every factor has common successors, so the ultraproduct does
too, by choosing coordinatewise successors. The affine criterion
for exactness therefore proves failure of exactness over every
nonzero unital ring. There is also no integral affine equalizer of this
pair, because reduction modulo two would give one. A total
category satisfies the integral pair condition IV, so this
ultraproduct is not total. The identical proof applies to the
ultrapower of $\Delta_{\rm inj}$.
\end{proof}

The preceding theorem concerns the ordinary category
ultraproduct, with all internal objects and arrows and with
ordinary finitely supported module elements. Replacing it by
a selected subcategory or by a construction permitting internal
hyperfinite sums changes the problem and is not covered by an
assertion of modular exactness or totality here.
